\section{Space, reachability, and finite counting}
\label{sec:space}

Space complexity makes the distinction between an arithmetical definition and
an efficient implementation particularly clear. A large finite table gives a
convenient primitive-recursive definition of reachability. A different,
depth-first implementation establishes a small space bound. The arithmetic
proof must connect the two. We give these connections explicitly, and then
reconstruct nondeterministic complementation and quantified Boolean
completeness in the same finite setting.

\subsection{Configurations and their bounds}

Fix a machine with $q$ states, $k$ work tapes, and a work alphabet of cardinality
$g\geq 2$. The binary input tape is read-only and has endmarkers. For an input
of length $n$, restrict each work tape to cells $0,\ldots,s-1$, with $s\geq 1$.
A configuration consists of the state, the input-head position, all work-head
positions, and all work-tape contents. Consequently
\[
 N_0=q(n+2)s^k g^{ks}
\]
is an exact bound on syntactic configurations. Mixed-radix encoding gives
primitive-recursive encoding and projection functions. In particular, tape
contents alone would give an incorrect configuration count: head positions
and the state matter.

The configuration graph contains an edge precisely when a legal machine
transition connects the two configurations without exceeding the cap. Delete
moves outside the cap. Add one new target vertex, numbered $N_0$, replace
the outgoing edges of every accepting configuration by its sole edge to
that target, and give the target no outgoing edges. This preserves
acceptance and keeps deterministic configuration graphs deterministic.
Put $N=N_0+1$. Initial and
target vertices are explicit functions of the machine and input. In a padded
binary encoding, unused codes simply have no useful outgoing edges.

Let $\ell$ be the least integer with $2^\ell\geq N+1$. For a fixed machine,
\[
 \ell\leq C_M\bigl(s+\lceil\log_2(n+2)\rceil+1\bigr).
\]
Here and below logarithm notation abbreviates integer binary length, defined
by bounded search among powers of two. The bound follows from the product
formula and $\log_2 s\leq s$ for $s\geq1$; it involves no analysis of real
logarithms.

Acceptance in the capped graph means existence of an accepting branch that
stays inside the cap. This agrees with ordinary acceptance when all branches
of the original machine obey the stated space bound. Branches may loop. If a
physical simulation must compute a function $s(n)$, that computation must
also obey the claimed space bound. The logarithmic and polynomial caps used
below have explicit suitable constructors. We make no unrestricted claim
that computing an arbitrary primitive-recursive cap is inexpensive.

\subsection{Paths, cycle removal, and the deterministic cutoff}

A path of length $m$ is a coded sequence $v_0,\ldots,v_m$ of vertices below
$N$, with the required adjacent edges. Its code is bounded by $N^{m+1}$,
with the length supplied separately. The predicate
$\mathrm{PathLe}(a,b,m)$ is a bounded disjunction over lengths $r\leq m$ and
codes below $N^{r+1}$. This deliberately expensive definition immediately
establishes primitive recursiveness.

\begin{lemma}[Constructive cycle removal]
Given a path from $a$ to $b$ in an $N$-vertex graph, a primitive-recursive
operation produces a path with the same endpoints and length at most $N-1$.
\end{lemma}
\begin{proof}
If positions $i<j$ contain the same vertex, delete positions
$i+1,\ldots,j$. The edge after position $i$ remains legal because the old
vertices at $i$ and $j$ coincide. Endpoints are unchanged, and the length
drops by at least one. Find the least repeated pair by bounded search, and
perform this deletion at most $m$ times, doing nothing once there is no
repeat. This is a fixed bounded iteration. Its invariant says that endpoints
and adjacency are preserved and that either no repeat remains or the length
has decreased. After $m$ rounds no repeat can remain. A duplicate-free list
of vertices below $N$ has at most $N$ entries, by finite counting.
\end{proof}

All assertions in this proof are bounded properties of the original path and
its successive codes. Their characteristic functions turn the invariant into
a quantifier-free formula. The finite counting assertion can itself be proved
using indicator sums over the $N$ possible vertex labels.

For a deterministic capped machine, an accepting trace can therefore be
shortened to one with fewer than $N$ edges. Determinism ensures that this
short trace agrees with the initial segment of the unique run. Thus the
following bounded implication holds for every supplied $t$: acceptance by
$t$ implies acceptance before $N$. A capped recognizer can be made a decider
for its acceptance language by rejecting at this cutoff. There is no need to
prove an unbounded statement about every future configuration as an
induction invariant.

\subsection{A polynomial graph algorithm}

For source $a$, let $S_0(v)=[v=a]$ and define
\[
 S_{i+1}(v)=S_i(v)\mathbin\lor
       \bigvee_{u<N}\bigl(S_i(u)\mathbin\land E(u,v)\bigr).
\]
Encode $S_i$ as an $N$-bit integer. Updating the vector is a fixed
primitive-recursive operation, so recursion on $i$ defines every vector.
An induction shows that $S_i(v)=1$ precisely when there is a path of length
at most $i$ from $a$ to $v$: a newly marked vertex is reached by extending
an old path, and the penultimate vertex of any nonempty path supplies the
converse. Witness paths can be stored in the mathematical construction,
using bounded least search, even when the implementation needs only bits.

The conjunction of this invariant over $v<N$ is bounded and can be
Boolean-compiled before induction. Cycle removal shows that $S_{N-1}$ is
the full reachable set. Computing $N$ rounds by scanning all ordered pairs
uses $O(N^3)$ adjacency queries and polynomial time on an ordinary explicit
graph representation. For a logarithmic-space machine the configuration
count is polynomial in input length, and adjacency is polynomial-time
decidable. This gives an explicit proof of $\mathrm{NL}\subseteq\mathrm{P}$.

\subsection{Savitch's theorem with a primitive-recursive audit}

The classical square-space simulation is due to Savitch
\cite{Savitch1970}. Define Boolean tables
\begin{align*}
 H_0(a,b)&=[a=b]\lor E(a,b),\\
 H_{j+1}(a,b)&=\bigvee_{z<N}
               \bigl(H_j(a,z)\land H_j(z,b)\bigr).
\end{align*}
For a direct primitive-recursive definition, compute the entire $N$ by $N$
table and update it by this formula. Table codes are below $2^{N^2}$.
Recursion on a single table code is ordinary primitive recursion; an
unanalysed recursive subroutine is not being treated as a totality axiom.

\begin{theorem}[Finite repeated squaring]
For every $j$, $H_j(a,b)=1$ if and only if there is a path from $a$ to $b$
of length at most $2^j$.
\end{theorem}
\begin{proof}
The base case permits lengths zero and one. Concatenating two paths of
length at most $2^j$ gives a path of length at most $2^{j+1}$. Conversely,
in a path of length $m\leq2^{j+1}$, take the vertex at position
$\min(m,2^j)$. The two resulting segments both have length at most $2^j$.
This midpoint witnesses the displayed disjunction. Split and join operations
on the coded paths are primitive recursive.
\end{proof}

For the iPRA induction, bound path codes by $N^{2^j+1}$ and conjoin over
$a,b<N$. The resulting assertion is bounded, even though its bound is large.
Its characteristic function supplies the quantifier-free invariant.
If $j$ is the least integer with $2^j\geq N$, repeated squaring decides
reachability.

The small-space implementation evaluates this recurrence depth first. Each
frame stores $a,b$, the midpoint loop index, the current level, a phase flag,
and one child answer. Evaluate one child, retain only its Boolean result,
then evaluate the other. At most $j+1$ frames are active. Each frame uses
$O(\ell+\log(j+2))$ bits, and $j\leq\ell$, giving $O(\ell^2)$ bits in all,
plus reusable adjacency scratch space. The recursion tree has at most
\[
 1+(2N)+\cdots+(2N)^j\leq(2N+1)^{j+1}
\]
nodes. Multiplication by the explicit per-node bookkeeping and adjacency
cost gives a primitive-recursive clock for an iterative stack machine.

Quantifier-free induction on the level, bounded over endpoints and the
relevant stack codes, proves that this clocked evaluator returns the table
value and obeys its frame bound. This is the essential bridge between the
large table definition and the small-space implementation. The table alone
would not prove a space bound.

\begin{theorem}[Square-space simulation for explicit caps]
For the explicit machine and cap data above, nondeterministic acceptance is
decidable in deterministic space
$O((s+\log(n+2))^2)$. If $s\geq\lceil\log_2(n+2)\rceil$, the bound is
$O(s^2)$. In particular, $\mathrm{NPSPACE}=\mathrm{PSPACE}$.
\end{theorem}
\begin{proof}
Use the configuration graph and the depth-first evaluator. The vertex code
length is $O(s+\log(n+2))$. Testing one transition needs only comparable
scratch space. For a polynomial cap its square is polynomial; the reverse
containment uses a deterministic machine as a nondeterministic machine with
one possible transition.
\end{proof}

\subsection{Nondeterministic complementation by exact counting}

The counting method behind nondeterministic complementation is due
independently to Immerman and Szelepcs\'enyi; we cite Immerman's account
\cite{Immerman1988}. Here we give the finite certificate construction and
its restricted-induction audit.

Fix source $a<N$ and put
\[
 R_i(v)=[\mathrm{PathLe}(a,v,i)],\qquad
 c_i=\sum_{v<N}R_i(v).
\]
In particular $c_0=1$. We will compute $c_{N-1}$ nondeterministically so
that every surviving branch returns the correct count and at least one
branch survives.

\begin{lemma}[Exact enumeration]
Suppose the correct value $c_i$ is supplied. A nondeterministic scan using
$O(\log N)$ work bits, plus adjacency scratch space, can enumerate precisely
the vertices reachable within $i$ steps. Every successful scan is exact,
and some scan succeeds.
\end{lemma}
\begin{proof}
Visit vertices $u=0,\ldots,N-1$ in order and guess a mark $b_u$. If $b_u=1$,
guess a path from $a$ to $u$ of length at most $i$, checking one edge at a
time with a bounded step counter. Reject this branch if the path is invalid.
Count all marks and require that their total equal $c_i$.

On a successful branch, $b_u\leq R_i(u)$ for every $u$. Equality of their
finite sums forces equality at every vertex: omitting one reachable vertex
would reduce the sum by at least one. Conversely mark exactly the reachable
vertices and use, for each, the least valid bounded path code. These choices
give a successful branch. Their mathematical constructor is bounded search,
so it is primitive recursive; it need not be a logarithmic-space algorithm.
The verifier stores only a constant number of vertices and counters, never
all marks or paths.
\end{proof}

Given $c_i$, compute $c_{i+1}$ as follows. For each target $v<N$, run a fresh
exact enumeration and accumulate the bit
\[
 h_v=[v=a]\lor\bigvee_{u<N}\bigl(b_u\land E(u,v)\bigr).
\]
After the count check succeeds, this is exactly $R_{i+1}(v)$. Increment a
next-count register when it is one. The outer scan keeps only the old count,
the partial new count, and a constant number of indices. Every successful
branch produces $c_{i+1}$, and the canonical path choices show that a
successful branch exists. Iterate from $i=0$ through $N-2$.

Finally perform exact enumeration for depth $N-1$. Accept precisely when
the enumeration succeeds and the target $b$ was not marked. Cycle removal
then gives
\[
 \text{an accepting branch exists}
 \quad\Longleftrightarrow\quad
 b\text{ is not reachable from }a.
\]
This construction does not reverse nondeterministic accept and reject
states. Its complement certificate establishes an exact count and uses
that count to exclude the target.

There are $O(N^4)$ adjacency queries: at most $N$ levels, $N$ targets per
level, $N$ enumerated vertices per target, and at most $N$ edges per guessed
path. Bookkeeping has polynomial cost. Space is $O(\log N)$ plus one
adjacency check. Thus the explicit logarithmic cap construction yields
$\mathrm{NL}=\mathrm{coNL}$. Applied to polynomial-space configuration
graphs, it also gives nondeterministic polynomial-space complementation.

\paragraph{Why the induction is permitted.}
A transcript has a primitive-recursive bit-length bound because all four
loops and every path are bounded. For example, a padded format has at most
$N(N-1)+1$ enumerations; each contains $N$ mark entries, each with a length
field and at most $N$ vertex fields. With
$w=\lceil\log_2(N+2)\rceil$, its bit length is bounded by
\[
 B(N)=\bigl(N(N-1)+1\bigr)N(1+w+Nw).
\]
Use the corresponding code bound $2^{B(N)}$ for quantified transcript
variables. The stage invariant says that every surviving bounded transcript
has the correct count, and that a specified primitive-recursive canonical
transcript survives. The first assertion is bounded universal; the second
is a Boolean equality after the canonical constructor has been defined.
Both therefore compile to quantifier-free induction. The canonical
constructor uses the finite reachability table and bounded least path
search. No unbounded choice or collection is needed. A formal proof must
also show that the transcript verifier and the low-space machine describe
exactly the same bounded branches; this is a local finite realization lemma.

\subsection{True quantified Boolean formulas are PSPACE-complete}
\label{sec:tqbf}

A quantified Boolean formula has constants, Boolean connectives, and
quantifiers over the two-element domain. Its truth is a primitive-recursive
Boolean function of its finite syntax code: bounded syntax traversal and
finite assignment tables provide a direct definition. A depth-first
evaluator instead stores the current assignment and a stack of subformula
positions. For an input of length $L$ there are at most $L$ active frames,
with polynomially many bits in total. Its finite recursion tree has an
explicit exponential bound, so it can be implemented with a clock. Thus
truth of closed quantified Boolean formulas lies in PSPACE.

For hardness, consider a fixed polynomial-space machine on input $x$.
Choose fixed-width binary configuration fields. Let $d$ bits suffice for a
configuration, including a unique accepting target. Validity and adjacency
are computable by Boolean circuits of size polynomial in $|x|+s(|x|)$:
compare states and heads, check unchanged tape cells, inspect the selected
work symbols, and select the input bit by its head index. Each comparison,
selector, and local transition has an explicit finite gate construction.
There are at most $2^d$ possible codes. Invalid codes have no useful edges.

The equality-or-adjacency predicate has a polynomial-size Boolean formula
with existentially quantified auxiliary gate variables. Introduce one
variable per circuit gate, constrain it to equal that gate's Boolean
operation, and require the output bit to be one. This representation is
linear in circuit size; expanding a circuit into a formula tree would not
justify a polynomial bound. Gate values supply canonical witnesses, and
induction in gate order proves the converse.

Call this base formula $\Phi_0(a,b)$, with $a,b$ tuples of $d$ free Boolean
variables, and require endpoint validity. Recursively define
\[
\begin{split}
 \Phi_{j+1}(a,b)=\exists z\bigl(&\mathrm{Valid}(z)\land
    \forall u\,\forall v\,\bigl(
       G(a,b,z,u,v)\to\Phi_j(u,v)\bigr)\bigr),\\
 G(a,b,z,u,v)={}&(u=a\land v=z)\lor(u=z\land v=b).
\end{split}
\]
Each displayed quantifier ranges over a $d$-bit tuple, and every level uses
fresh variable names. There is only one recursive occurrence of $\Phi_j$.
Validity can be included with a polynomial-size circuit and fresh
existential gate variables at each level; its contribution remains
polynomial. If all $d$-bit codes are permitted vertices and invalid codes
are isolated, midpoint validity may instead be omitted: a path between
valid endpoints cannot benefit from an isolated invalid midpoint.

\begin{lemma}[Shared occurrence of reachability]
For valid endpoints, $\Phi_j(a,b)$ is true exactly when there is a path of
length at most $2^j$ from $a$ to $b$.
\end{lemma}
\begin{proof}
The base formula has the equality-or-adjacency meaning. If the two half-paths
exist for midpoint $z$, the induction hypothesis proves the consequent for
each of the two guarded endpoint pairs. For any other $u,v$ the Boolean
guard is false, and its implication is true. This proves the universal
formula. Conversely instantiate its universal tuples first by $(a,z)$ and
then by $(z,b)$. Both guards are true, so the two instances of the induction
hypothesis give the half-paths. Concatenate them. Every case distinction
is on a decidable finite Boolean equality.
\end{proof}

Take $j=d$. Cycle removal makes this depth sufficient, and substitute the
initial and accepting configurations. The resulting closed formula is true
exactly when the original machine accepts. The guard and quantified tuples
add $O(d)$ variable occurrences per level. The base circuit and any repeated
validity circuits have polynomial size; with $d$ levels the full formula
has polynomial size in $|x|+s(|x|)$. Binary variable identifiers add only a
logarithmic factor. A printer performs explicitly bounded nested loops over
the fields, gates, levels, and variable indices, giving a polynomial-time
many-one reduction for each fixed polynomial-space machine and its cap.

Quantifiers need not initially be in prenex form. If that convention is
required, rename bound variables, put the Boolean matrix in negation normal
form, and move each quantifier across a connective where its variable is
not free in the other operand. These finite Boolean equivalences preserve
truth and require no duplication of subformulas, giving polynomial output
size. The domain is nonempty, and all Boolean subformulas are decidable.

The iPRA invariant quantifies only over endpoint tuples below $2^d$, finite
syntax codes under the explicit size bound, and assignment codes. Truth
and repeated-squaring values are primitive-recursive Booleans. Their
equality, conjoined over bounded endpoints, is therefore an admissible
quantifier-free induction invariant after compilation. This proves the
arithmetical content of TQBF PSPACE-completeness as a family of explicit
reduction and correctness statements, without quantifying over arbitrary
sets of inputs.

\subsection{Containments and implementation obligations}

The constructions give
\[
 \mathrm{L}\subseteq\mathrm{NL}=\mathrm{coNL}
 \subseteq\mathrm{P}\subseteq\mathrm{NP}
 \subseteq\mathrm{PSPACE}=\mathrm{NPSPACE}
 \subseteq\mathrm{EXPTIME}.
\]
Deterministic computations supply single-branch nondeterministic ones.
For $\mathrm{NP}\subseteq\mathrm{PSPACE}$, enumerate all polynomial-length
witnesses with a binary counter, rerun the polynomial-time verifier, and
reuse its storage. The counter and verifier need polynomial space although
the total time is exponential. For
$\mathrm{PSPACE}\subseteq\mathrm{EXPTIME}$, the deterministic configuration
cutoff is $2^{p(n)}$ for an explicit polynomial $p$ depending on the fixed
machine. Simulating that many transitions with polynomial overhead preserves
an exponential-time bound. None of these arguments decides any of the
standard unresolved equalities between neighboring classes.

For Lean, the semantic layer may start with finite vertices and a Boolean
adjacency function. The arithmetic layer must replace this function by a
finite table or a particular coded transition computation. Required bridge
lemmas include the configuration encoding bound, transition--edge
correspondence, cycle-removal correctness, repeated-squaring correctness,
stack-evaluator equivalence and frame bound, the equal-count finite subset
lemma, counting-transcript soundness and completeness, branch realization,
and the shared-occurrence formula lemma. Finally the bounded invariants
must be translated into the chosen object-language proof calculus.

The executable reference implementation supplied with this report checks
these finite constructions on small graphs. Its transcript objects and
independent closure oracle are intentionally materialized; the Python code
is not claimed to be a logarithmic-space Turing-machine implementation.
Finite tests validate examples and detect implementation errors. They do
not replace the mathematical arguments or an eventual checked iPRA
derivation.
